\section{Local zeta values and consequences of sparse fourth moments}\label{app:localzeta}

\begin{proposition}[local values from harmonic measure]\label{prop:localzetaproof}
For every fixed $p>0$,
\begin{equation}\label{eq:baselocal}
 |\zeta(1/2+it)|^p\ll_p
 1+\log(2+|t|)\int_{t-1/3}^{t+1/3}|\zeta(1/2+iu)|^p\,du.
\end{equation}
\end{proposition}
\begin{proof}
We give the localization argument, including the analytic estimates it
requires. The elementary functional equation gives, uniformly for
$0\le u\le1/\log T$ and $\tau\asymp T$,
\begin{equation}\label{eq:localreflection}
 |\zeta(1/2-u+i\tau)|\ll |\zeta(1/2+u+i\tau)|.
\end{equation}
Here is one derivation of the facts needed in this statement. Periodizing
the Gaussian, whose Fourier transform is $x^{-1/2}e^{-\pi\xi^2/x}$,
gives $\theta(x)=x^{-1/2}\theta(1/x)$. Mellin integration, initially
for $\Re s>1$, and splitting at $x=1$, then gives
\[
 \pi^{-s/2}\Gamma(s/2)\zeta(s)
 =-\frac1s-\frac1{1-s}
 +\frac12\int_1^\infty(\theta(x)-1)
       \bigl(x^{s/2}+x^{(1-s)/2}\bigr)\,\frac{dx}{x}.
\]
The right side continues meromorphically and is unchanged by
$s\mapsto1-s$. The gamma quotient therefore has modulus
$\asymp(1+|\tau|)^u$ at $s=1/2-u+i\tau$, which proves
\eqref{eq:localreflection}. The gamma estimate follows
by applying Euler summation twice to the logarithmic gamma product:
on $\Re z>0$ it gives
\[
 \log\Gamma(z)=(z-\tfrac12)\log z-z+c+\frac1{12z}
 -\frac12\int_0^\infty\frac{B_2(\{x\})}{(x+z)^2}\,dx.
\]
The integral is $O(|\Im z|^{-1})$ when $\Re z$ is in a fixed positive
compact interval; take real parts in the two gamma factors.
Conjugation accounts for the sign of $\tau$. Also Euler summation in
the defining Dirichlet series, with a cutoff $X\asymp1+|\tau|$, gives
\[
 \zeta(s)=\sum_{n\le X}n^{-s}+\frac{X^{1-s}}{s-1}
          +O(X^{-\Re s}+|s|X^{-\Re s}/\Re s)
\]
for $\Re s>0$ away from $s=1$. It gives the polynomial growth used
below directly from Euler summation.

We prove the bound for $T\le t\le2T$ with $T\ge10$. On bounded intervals, continuity gives the result.
Set $s_0=1/2+it$, $\delta=1/\log T$, and choose a fixed sufficiently
large $C=C_p$. On the strip $1/2\le\Re s\le3/2$ put
\[
 G(s)=\frac{s-1}{s+1}\zeta(s)
 \exp\bigl(C\log T\bigl[(s-s_0)^2-(s-s_0)\bigr]\bigr).
\]
The pole is canceled. On either boundary the exponential has modulus
$\exp(-C\log T\,|\Im s-t|^2)$. On the right boundary $G$ is bounded
by an absolute constant. In the small rectangle
$0\le\Re(s-s_0)\le\delta$, $|\Im(s-s_0)|\le\delta$,
the other factors show $|\zeta(s)|\ll_C|G(s)|$.

The harmonic-measure kernel for the left boundary of the strip of width
one is
\[
 K_u(v)=\frac{\sin(\pi u)}{2(\cosh(\pi v)-\cos(\pi u))},
 \qquad 0<u<1.
\]
This formula follows by mapping $0<\Re z<1$ to the upper half-plane
by $e^{\pi i z}$ and changing variables in the Poisson kernel.
It is positive, has integral $1-u\le1$, and is bounded by
$C u/(u^2+v^2)$ for $u\le1/2$. Subharmonicity of $|G|^p$ and the
maximum principle on finite rectangles, followed by a limit, give
\[
 |\zeta(1/2+u+i\tau)|^p
 \ll_{p,C}1+\int_{\mathbb R}K_u(\tau-v)
 |\zeta(1/2+iv)|^p T^{-pC(v-t)^2}\,dv
\]
in the small rectangle. The horizontal boundary terms vanish because
of Gaussian decay and the polynomial bound just obtained.

Apply the disk submean inequality to $|\zeta|^p$ at $s_0$ with radius
$\delta$. Reflect the left half using \eqref{eq:localreflection},
and enlarge the right half to the small rectangle. For every $v$,
\[
 \int_0^\delta\int_{t-\delta}^{t+\delta}K_u(\tau-v)\,d\tau\,du
 \le\delta.
\]
After division by the disk area, the bound is
\[
 |\zeta(s_0)|^p\ll_p
 1+\delta^{-1}\int_{\mathbb R}
 |\zeta(1/2+iv)|^pT^{-pC(v-t)^2}\,dv.
\]
For $|v-t|>1/3$, the polynomial bound and sufficiently large $C_p$
make this integral $O(\delta)$, including $|v|$ far larger than $T$.
For the remaining interval discard the Gaussian. This proves
\eqref{eq:baselocal}. The disk submean inequality itself follows from
the circle mean inequality for the subharmonic function $|\zeta|^p$
by integrating the radius, including disks containing zeros.
\end{proof}

\begin{proposition}[fourth and twelfth moments from one sparse input]\label{prop:momentbaseproof}
Assuming only the sparse-fourth-moment estimate
\eqref{eq:basefourth} among the deep zeta moment inputs, one has
\[
 M_4(T)\prec T,\qquad M_{12}(T)\prec T^2,\qquad
 \cN(T,V)\prec\min(TV^{-4},T^2V^{-12})\quad(V\ge1).
\]
\end{proposition}
This derivation obtains the classical twelfth moment from the
sparse fourth moment. Attribution remains with the original authors;
this implication is Iwaniec's route~\cite{Iwa80}.
\begin{proof}
Take $J=1,L=T$ in \eqref{eq:basefourth} and sum dyadic height intervals
to obtain $M_4(T)\prec T$. Proposition~\ref{prop:localzetaproof}
with $p=4$ and disjoint neighborhoods of a $1$-separated set gives
$RV^4\prec T$; the added constant is absorbed for $V$ sufficiently
large, and the bounded-$V$ range is trivial.

Fix $\rho>0$, with all occurrences of $\prec$ in this proof taken at
a loss much smaller than $\rho$. If $V^4\le T^{1/2+\rho}$, the fourth
moment gives
$R\prec TV^{-4}\le T^{2+2\rho}V^{-12}$.
Otherwise put $L=cV^4T^{-\rho}$, with a fixed positive constant $c$.
If necessary increase the lower threshold by a constant so that
$L\ge\sqrt T$; the transition is covered by the preceding case.
When $R\ge1$, the just-proved fourth-moment estimate and the local
inequality give $V^4\ll T^{1+\rho/2}$, hence $L<T$ for large $T$.

Cover the disjoint neighborhoods $[t_j-1/3,t_j+1/3]$ by grid intervals
of length $L$. At most $2R$ grid intervals are needed. Boundary
intervals at $T,2T$ may be enlarged or shifted; a fixed number of
families in comparable dyadic height ranges handles them. Thus
\eqref{eq:basefourth} and the local lower mass yield
\[
 RV^4\prec RL+T\sqrt{R/L}.
\]
The $RL$ term is absorbed, since $L\ll V^4T^{-\rho}$.
Squaring the remaining inequality gives
$R\prec T^2/(LV^8)\ll T^{2+\rho}V^{-12}$.
As $\rho$ is arbitrary this proves the discrete twelfth-moment bound.

Finally the measure of each amplitude layer is at most a constant times
the maximal $1$-separated cardinality: choose a maximal separated
subset, whose unit neighborhoods cover that layer. Integrate in
dyadic $V$ and use the polynomial pointwise bound from the preceding
proof to ensure only $O(\log T)$ layers. This gives $M_{12}(T)\prec T^2$.
\end{proof}

\section{The derivative estimate from a precise mean-value input}\label{app:hbderivative}

We rederive the Heath-Brown estimate used here, keeping its attribution
to \cite{HB17}. Its sole nonelementary input in this appendix is
\begin{equation}\label{eq:vmvtinput}
 J_{s,q}(H):=\int_{[0,1]^q}\left|\sum_{h=1}^{\lfloor H\rfloor}
 e(\alpha_1h+\cdots+\alpha_qh^q)\right|^{2s}d\boldsymbol\alpha
 \ll_{q,\epsilon}H^{s+\epsilon},\qquad s=q(q+1)/2.
\end{equation}
This is the optimal Vinogradov mean-value theorem \cite{BDG16,Woo16}.

\begin{lemma}[coefficient clustering]\label{lem:hbcluster}
Suppose $d\ge3$ and $\lambda\le f^{(d)}\le A\lambda$ on an interval
of length $M$, with $A$ fixed. Put $H=\lfloor(A\lambda)^{-1/d}\rfloor$.
If $A\lambda$ is a sufficiently small constant depending on $d$, then
\[
 \mathcal P:=\#\left\{m,n:\left\|\frac{f^{(j)}(m)-f^{(j)}(n)}{j!}\right\|
 \le2H^{-j}\ (1\le j<d)\right\}
 \ll_{A,d}(M+\lambda M^2+\lambda^{-2/d})\log(2M).
\]
\end{lemma}
\begin{proof}
Keep the last two derivative conditions to obtain an upper bound. Let
$g_1=f^{(d-2)}/(d-2)!$, $g_2=f^{(d-1)}/(d-1)!$, and periodize
triangular functions of widths $B=4H^{2-d}$ and $C=4H^{1-d}$.
Their product $P(x,y)$ is at least $1/4$ on the counted differences,
vanishes unless $\|x\|\le B,\|y\|\le C$, and has nonnegative Fourier
coefficients $BC\,\operatorname{sinc}(rB)^2\operatorname{sinc}(sC)^2$.
Smallness of $\lambda$ ensures $B,C<1/2$.

Partition the integer interval into $K\asymp1+A\lambda M$ consecutive
pieces whose lengths make $|g_2(m)-g_2(n)|<1/2$ within each piece.
Fourier expansion and $|\sum_{i\le K}z_i|^2\le K\sum_i|z_i|^2$ show
\[
 \mathcal P\ll K\sum_i\sum_{m,n\text{ in piece }i}
 P(g_1(m)-g_1(n),g_2(m)-g_2(n)).
\]
Within a piece the $g_2$ condition is an ordinary real difference,
and monotonicity gives $|m-n|\ll\lambda^{-1}H^{1-d}\asymp H$.
The diagonal costs $M$. For $m=n+h$, $1\le h\ll H$, the derivative
of $g_1(x+h)-g_1(x)$ lies between fixed positive multiples of
$\lambda h$. A monotone function with derivative in $[\mu,A'\mu]$
takes integer arguments within distance $b$ of an integer at most
$O((1+\mu M)(1+b/\mu))$ times: count the integers in its range and
the length of each inverse interval. Hence the off-diagonal
count, even after dropping the piece restriction, is at most
\[
 \sum_{h\ll H}(1+\lambda Mh)(1+B/(\lambda h))
 \ll H+\lambda MH^2+(B/\lambda)\log(2H)+MBH
 \ll M+H^2\log(2M).
\]
Here $\lambda\asymp H^{-d}$ and $d\ge3$; if $H>M$ the trivial
$M^2\le H^2$ bound suffices. Multiply by $K$ and use
$\lambda M H^2\ll M$. This proves the bound.
\end{proof}

\begin{proposition}[Heath-Brown from \eqref{eq:vmvtinput}]
\label{prop:hbderivativeproof}
The mean-value bound \eqref{eq:vmvtinput} with $q=d-1$ implies
\eqref{eq:basederivative} for $d\ge3$.
\end{proposition}
\begin{proof}
Write $s=d(d-1)/2$ and first consider an integer interval of length $M$.
Conjugation makes $f^{(d)}$ positive. If $\lambda$ is bounded below,
the claimed first term bounds the sum trivially. If $H\ge M$, the
third claimed term does so. Otherwise $1\ll H<M$ and averaging
the $H$ translates of the sum, with end intervals costing $O(H)$, gives
\[
 \left|\sum_ne(f(n))\right|
 \ll H+H^{-1}\sum_n\max_{x\le H}
 \left|\sum_{h\le x}e\left(\sum_{j=1}^{d-1}\frac{f^{(j)}(n)}{j!}h^j\right)\right|.
\]
Taylor's remainder was removed by Abel summation: its derivative is
$O_d(A\lambda H^{d-1})=O_d(H^{-1})$.
To avoid a point-dependent maximum, use Fourier completion on
$\mathbb Z/(2H)\mathbb Z$, as in Lemma~\ref{lem:completionproof}.
This costs $O(\log(2H))$ and adds a single common linear coefficient
to every polynomial, preserving all coefficient differences.

Let $\nu(\boldsymbol\alpha)$ count the $n$ for which all the polynomial
coefficients are within $H^{-j}$ of $\alpha_j$ modulo one. Perturbing
coefficients by these amounts has total phase variation $O_d(1)$
on $[0,H]$. Abel summation and averaging over the coefficient box
therefore bound the preceding $n$ sum by
\[
 C_dH^s\left(I(H)+H^{-1}\int_0^H I(x)\,dx\right),\qquad
 I(x)=\int_{[0,1]^{d-1}}\nu(\boldsymbol\alpha)
 \left|\sum_{h\le x}e(\alpha_1h+\cdots+\alpha_{d-1}h^{d-1})\right|
 d\boldsymbol\alpha.
\]
Box volumes and overlaps give
$\int\nu\ll MH^{-s}$ and $\int\nu^2\ll\mathcal P H^{-s}$.
H\"older with three factors, of reciprocal exponents
$1-1/s$, $1/(2s)$ and $1/(2s)$, gives
\[
 I(x)\le(\int\nu)^{1-1/s}(\int\nu^2)^{1/(2s)}
 J_{s,d-1}(H)^{1/(2s)}.
\]
The monotonicity in $H$ follows directly by expanding the mean value
as a number of solutions. Applying \eqref{eq:vmvtinput} and collecting
the powers of $H$, which cancel, yields
\[
 \left|\sum_ne(f(n))\right|
 \ll_{d,A,\epsilon}H+M^{1-1/s+\epsilon}\mathcal P^{1/(2s)}.
\]
Insert Lemma~\ref{lem:hbcluster}. Since $2s=d(d-1)$, this is
\[
 \ll M^\epsilon\left(\lambda^{-1/d}
 +M^{1-1/[d(d-1)]}+M\lambda^{1/[d(d-1)]}
 +M^{1-2/[d(d-1)]}\lambda^{-2/[d^2(d-1)]}\right).
\]
If $\lambda<M^{-d}$ the last term is at least $M$. Otherwise
$\lambda^{-1/d}\le M$, and the last displayed term bounds it. It can therefore be removed in both cases. Substitute
$\lambda\asymp FM^{-d}$ to obtain \eqref{eq:basederivative}.
For a shorter subinterval first use its own length in this argument;
each term is increasing in that length at fixed $\lambda$, so replacing
it by $M$ is valid. The argument uses the phase on its given interval.
\end{proof}
